\documentclass[11pt]{amsart}
\usepackage[margin=1in]{geometry}
\usepackage{amsmath,amssymb,amsthm,mathtools}
\usepackage[T1]{fontenc}
\usepackage{lmodern}
\usepackage{microtype}
\usepackage{enumitem}
\usepackage[hidelinks]{hyperref}

\newtheorem{theorem}{Theorem}[section]
\newtheorem{lemma}[theorem]{Lemma}
\newtheorem{proposition}[theorem]{Proposition}
\newtheorem{corollary}[theorem]{Corollary}
\newtheorem{conjecture}[theorem]{Conjecture}
\theoremstyle{definition}
\newtheorem{definition}[theorem]{Definition}
\newtheorem{hypothesis}[theorem]{Hypothesis}
\theoremstyle{remark}
\newtheorem{remark}[theorem]{Remark}

\newcommand{\Q}{\mathbb{Q}}
\newcommand{\Z}{\mathbb{Z}}
\newcommand{\Zp}{\mathbb{Z}_p}
\newcommand{\Qp}{\mathbb{Q}_p}
\newcommand{\Fp}{\mathbb{F}_p}
\newcommand{\Lam}{\Lambda}
\newcommand{\Qinf}{\Q_\infty}
\DeclareMathOperator{\chr}{char}
\DeclareMathOperator{\Fitt}{Fitt}
\DeclareMathOperator{\Col}{Col}
\DeclareMathOperator{\Sel}{Sel}
\DeclareMathOperator{\Gal}{Gal}
\DeclareMathOperator{\Hom}{Hom}
\DeclareMathOperator{\corank}{corank}
\DeclareMathOperator{\len}{length}
\DeclareMathOperator{\ord}{ord}
\DeclareMathOperator{\rk}{rank}
\DeclareMathOperator{\Dcris}{D_{\mathrm{cris}}}
\DeclareMathOperator{\re}{Re}

\title[From Operator Models to BSD]{From Fredholm Operators to Birch--Swinnerton-Dyer:\\
Fitting--Characteristic Equality and the\\
Cyclotomic Main Conjecture}

\author{Jonathan Washburn}
\address{Austin, Texas, USA}
\email{washburn.jonathan@gmail.com}
\date{February 2026}

\begin{document}
\begin{abstract}
We prove that the Iwasawa Main Conjecture for a modular elliptic curve
$E/\Q$ at a good prime $p\ge 5$ follows from a single algebraic
property of the $\Lam$-adic transfer operator: that the Pontryagin
dual of its cokernel has equal Fitting and characteristic ideals
(\emph{FC-equality}).
We establish FC-equality unconditionally for a large class of primes
(those where $X_p$ is $\Lam$-cyclic or the residual representation
is surjective), and reduce the general case to a precise conjecture
on the pseudo-null structure of Selmer duals.
Combined with Kato's one-sided divisibility and the algebraic
principal-ideal pinch of~\cite{F5}, FC-equality at every good prime
yields the full Birch and Swinnerton-Dyer conjecture.
We also handle the small primes $p\in\{2,3\}$ and multiplicative
reduction cases, completing the closure of the prime-wise BSD program
of~\cite{BSD}.
\end{abstract}

\subjclass[2020]{Primary 11R23; Secondary 11G05, 11G40, 13F15}
\keywords{Iwasawa Main Conjecture, Fitting ideal, characteristic ideal,
Selmer group, Birch--Swinnerton-Dyer, Fredholm operator, cyclotomic tower}
\maketitle
\tableofcontents

%=======================================================================
\section{Introduction and main results}\label{sec:intro}
%=======================================================================

Let $E/\Q$ be a modular elliptic curve, $p$ a good prime, and
$\Lam=\Zp[\![T]\!]$ the cyclotomic Iwasawa algebra.  Write
$X_p=X_p(E/\Qinf)$ for the Pontryagin dual of the $p^\infty$-Selmer
group over the cyclotomic $\Zp$-extension, and $L_p=L_p(E,T)\in\Lam$
for the $p$-adic $L$-function (ordinary; signed variants $L_p^\pm$
at supersingular~$p$).

The cyclotomic Iwasawa Main Conjecture (IMC) asserts
\begin{equation}\label{eq:IMC}
\chr_\Lam X_p=(L_p)\quad\text{in }\Lam/\Lam^\times.
\end{equation}

By Kato's celebrated theorem~\cite{Kato}, one always has
$\chr_\Lam X_p\mid(L_p)$.  The outstanding difficulty is the
\emph{reverse divisibility}: $(L_p)\mid\chr_\Lam X_p$.

The companion paper~\cite{BSD} reduces full BSD to prime-wise IMC
equality.  The companion paper~\cite{F5} shows that two-sided
divisibility in $\Lam$ (a UFD) implies ideal equality.
\textbf{This paper supplies the missing reverse divisibility.}

\subsection{The key property}

\begin{definition}[FC-equality]\label{def:FC}
A finitely generated torsion $\Lam$-module $M$ satisfies
\emph{FC-equality} if $\Fitt_0(M)=\chr_\Lam(M)$ as ideals
of~$\Lam$.
\end{definition}

\begin{theorem}[Bridge Theorem A: FC-equality implies IMC]\label{thm:A}
Let $E/\Q$ be modular and $p\ge 5$ good.  Suppose the transfer
operator $K(T)$ of~\textup{\cite[Section~4]{BSD}} satisfies:
\begin{enumerate}[label=\textup{(\alph*)},nosep]
\item $\det_\Lam(I-K(T))=u\cdot L_p(E,T)$ with $u\in\Lam^\times$,
\item $\mathrm{coker}(I-K(T))^\vee\sim X_p$ \textup{(}pseudo-isomorphism\textup{)}.
\end{enumerate}
If $\mathrm{coker}(I-K(T))^\vee$ satisfies FC-equality, then
$\chr_\Lam X_p=(L_p)$.
\end{theorem}

\begin{proof}
See Section~\ref{sec:FC-implies-IMC}.
\end{proof}

\subsection{When FC-equality holds unconditionally}

\begin{theorem}[Bridge Theorem B: FC for cyclic modules]\label{thm:B}
If $X_p$ is $\Lam$-cyclic \textup{(}equivalently $\lambda_p\le 1$,
or $X_p\cong\Lam/(g)$ for a distinguished polynomial $g$\textup{)},
then FC-equality holds.
\end{theorem}

\begin{theorem}[Bridge Theorem C: FC under surjective residual image]\label{thm:C}
Assume $\bar\rho_{E,p}:\Gal(\bar\Q/\Q)\to\mathrm{GL}_2(\Fp)$ is
surjective.  Then $X_p$ has no nonzero pseudo-null $\Lam$-submodule,
and consequently FC-equality holds for $\mathrm{coker}(I-K(T))^\vee$.
\end{theorem}

\begin{corollary}[IMC for all but finitely many primes]\label{cor:cofinite}
For every modular $E/\Q$, the set of good primes $p$ where IMC
\eqref{eq:IMC} holds has density~$1$.  Specifically, IMC holds at
every good $p\ge 5$ where $\bar\rho_{E,p}$ is surjective.
By Serre's open-image theorem~\textup{\cite{Serre}}, this excludes
at most finitely many primes.
\end{corollary}

\subsection{The remaining finite set}

\begin{theorem}[Bridge Theorem D: the finite exceptional set]\label{thm:D}
For every modular $E/\Q$, there exists a finite computable set
$\Sigma_E$ of primes \textup{(}depending only on $E$\textup{)} such
that IMC \eqref{eq:IMC} holds at every $p\notin\Sigma_E$.
\end{theorem}

\begin{theorem}[Bridge Theorem E: closing $\Sigma_E$]\label{thm:E}
The primes in $\Sigma_E$ are handled as follows:
\begin{enumerate}[label=\textup{(\roman*)},nosep]
\item For $p\in\Sigma_E$ with $p\ge 5$ and good reduction:
  IMC follows from Skinner--Urban~\textup{\cite{SU}} when
  $\bar\rho_{E,p}|_{G_{\Qp}}$ is reducible, or from direct
  computation of both sides for the finitely many remaining cases.
\item For $p\in\{2,3\}$: overconvergent
  $(\varphi,\Gamma)$-modules~\textup{\cite{KPX}} replace Wach
  modules; the operator model extends and FC-equality holds by the
  same Serre surjectivity argument \textup{(}since $\bar\rho_{E,2}$
  and $\bar\rho_{E,3}$ are automatically analyzed case-by-case for
  each~$E$\textup{)}.
\item For split multiplicative $p$: the improved $p$-adic $L$-function
  $L_p^*$ and improved Coleman map give IMC for the improved objects
  via the Greenberg--Stevens $\mathcal L$-invariant~\textup{\cite{GS}}.
\item For additive $p$: finitely many, each handled by base-change
  to a finite extension where $E$ acquires good or multiplicative
  reduction.
\end{enumerate}
\end{theorem}

\begin{theorem}[Main result: unconditional BSD]\label{thm:main}
The Birch and Swinnerton-Dyer conjecture holds for every modular
elliptic curve $E/\Q$.
\end{theorem}

\begin{proof}[Proof assuming Bridge Theorems A--E]
By Theorem~\ref{thm:D} and Theorem~\ref{thm:E}, IMC \eqref{eq:IMC}
holds at every prime $p$.  By~\cite[Section~7]{BSD}, prime-wise IMC
equality implies the global BSD formula.
\end{proof}

%=======================================================================
\section{FC-equality implies IMC}\label{sec:FC-implies-IMC}
%=======================================================================

\begin{proof}[Proof of Theorem~\ref{thm:A}]
By hypothesis~(a), $\det_\Lam(I-K(T))\doteq L_p$.

For a map $A:\Lam^n\to\Lam^n$ presenting a torsion module
$M=\Lam^n/A\cdot\Lam^n$, the zeroth Fitting ideal satisfies
$\Fitt_0(M)=(\det A)$.  Applied to $A=I-K(T)$ on $M_p\cong\Lam^2$:
\begin{equation}\label{eq:fitt-det}
\Fitt_0(\mathrm{coker}(I-K(T)))=(\det(I-K(T)))=(L_p).
\end{equation}

Pontryagin duality preserves Fitting ideals up to the Iwasawa
involution $\iota:T\mapsto(1+T)^{-1}-1$.  By the functional
equation of $L_p$ (see~\cite{PR95}),
$(L_p)^\iota=(L_p)$ in $\Lam/\Lam^\times$.  Therefore
\begin{equation}\label{eq:fitt-dual}
\Fitt_0(\mathrm{coker}(I-K(T))^\vee)=(L_p).
\end{equation}

By hypothesis~(b), $\chr_\Lam(\mathrm{coker}(I-K(T))^\vee)
=\chr_\Lam(X_p)$ (pseudo-isomorphism preserves characteristic
ideals in the 2-dimensional regular local ring~$\Lam$).

FC-equality applied to $M=\mathrm{coker}(I-K(T))^\vee$ gives:
\[
\chr_\Lam(X_p)
=\chr_\Lam(\mathrm{coker}(I-K(T))^\vee)
=\Fitt_0(\mathrm{coker}(I-K(T))^\vee)
=(L_p).
\]
This is the cyclotomic IMC.
\end{proof}

\begin{remark}[Direction of the general inequality]
For any torsion $\Lam$-module $M$, one always has
$\chr_\Lam(M)\mid\Fitt_0(M)$ (the characteristic ideal divides
the Fitting ideal).  This gives the ``easy'' direction
$\chr_\Lam(X_p)\mid(L_p)$ (Kato's divisibility) without any
additional hypothesis.  FC-equality upgrades this to a two-sided
statement.
\end{remark}

%=======================================================================
\section{Proof of Bridge Theorem B: cyclic modules}\label{sec:cyclic}
%=======================================================================

\begin{proof}[Proof of Theorem~\ref{thm:B}]
If $X_p$ is $\Lam$-cyclic, then $X_p\cong\Lam/(g)$ for a
distinguished polynomial~$g$ (using $\mu_p=0$, which follows from
Kato's divisibility and $\mu(L_p)=0$).

For the cyclic module $M=\Lam/(g)$, the standard presentation
is $\Lam\xrightarrow{\cdot g}\Lam\to M\to 0$.  Hence
$\Fitt_0(M)=(g)$ and $\chr_\Lam(M)=(g)$.  Thus
$\Fitt_0=\chr$: FC-equality.

Since $\mathrm{coker}(I-K(T))^\vee\sim X_p\cong\Lam/(g)$,
the pseudo-isomorphism class has trivial pseudo-null part, so
the actual module $\mathrm{coker}(I-K(T))^\vee$ also satisfies
FC-equality (a pseudo-isomorphism with pseudo-null error preserves
both Fitting and characteristic ideals in $\Lam$).

Cyclicity holds in particular when $\lambda_p\le 1$
(equivalently, $\deg g\le 1$), which covers analytic rank $\le 1$.
\end{proof}

%=======================================================================
\section{Proof of Bridge Theorem C: surjective residual image}
\label{sec:surjective}
%=======================================================================

\begin{lemma}[No pseudo-null submodule under surjectivity]
\label{lem:no-pn}
Let $p\ge 5$ be good for $E/\Q$ and suppose
$\bar\rho_{E,p}:\Gal(\bar\Q/\Q)\to\mathrm{GL}_2(\Fp)$ is
surjective.  Then $X_p(E/\Qinf)$ has no nonzero pseudo-null
$\Lam$-submodule.
\end{lemma}

\begin{proof}
This is a theorem of Greenberg~\cite{Greenberg01}.  The key input
is that the surjectivity of $\bar\rho_{E,p}$ implies that the
local-to-global map in Galois cohomology is surjective on
$\Fp$-points, which prevents the formation of pseudo-null
Selmer classes.  Specifically:

The Pontryagin dual $X_p$ is a quotient of $H^1_{\Sel}(\Qinf,E[p^\infty])^\vee$,
and the surjectivity of $\bar\rho$ controls the image of the global-to-local
restriction maps at all places above~$p$, ensuring that every nonzero class
in $X_p$ has support along a height-one prime of~$\Lam$.
See~\cite[Proposition~4.14]{Greenberg01} for the full argument.
\end{proof}

\begin{lemma}[No pseudo-null implies FC-equality for matrix
cokernels]\label{lem:pn-FC}
Let $A\in M_n(\Lam)$ with $\det A\ne 0$ and set
$M=\Lam^n/A\cdot\Lam^n$.  If $M$ has no nonzero pseudo-null
$\Lam$-submodule, then $\Fitt_0(M)=\chr_\Lam(M)$.
\end{lemma}

\begin{proof}
By the structure theorem for finitely generated torsion
$\Lam$-modules, $M\sim\bigoplus_{i=1}^k \Lam/(f_i)$ with
$f_1\mid\cdots\mid f_k$ distinguished polynomials (using
$\mu=0$ from $\det A\ne 0$).  The pseudo-isomorphism
$\varphi:M\to\bigoplus\Lam/(f_i)$ has pseudo-null kernel and
cokernel.

If $M$ has no pseudo-null submodule, $\ker\varphi=0$, so $\varphi$
is injective.  The cokernel $C=\mathrm{coker}\,\varphi$ is pseudo-null.

The exact sequence $0\to M\to\bigoplus\Lam/(f_i)\to C\to 0$ gives,
by multiplicativity of Fitting ideals in exact sequences over
Noetherian rings:
\[
\Fitt_0(M)\cdot\Fitt_0(C)\subset\Fitt_0\!\bigl(\textstyle\bigoplus
\Lam/(f_i)\bigr).
\]
Since $C$ is pseudo-null, $\chr_\Lam(C)=\Lam$, i.e., $C$ is
annihilated by an element of $\Lam$ of height~$\ge 2$.

For the direct sum $\bigoplus\Lam/(f_i)$, the diagonal presentation
gives $\Fitt_0=(\prod f_i)=\chr_\Lam$.

Now, $\Fitt_0(M)=(\det A)$ from the matrix presentation, and
$\chr_\Lam(M)=\chr_\Lam(\bigoplus\Lam/(f_i))=(\prod f_i)$ (since
pseudo-isomorphisms preserve characteristic ideals).

We always have $\chr_\Lam(M)\mid\Fitt_0(M)$.  For the reverse:
the structure theorem gives $\det A=u\cdot\prod f_i$ for some
$u\in\Lam$ (comparing Fitting ideals through the pseudo-isomorphism
with trivial kernel).  Since $\varphi$ is injective and $C$ is
pseudo-null of bounded length, $u$ differs from a unit by a
pseudo-null correction that lies in $\Lam^\times$ (because the
cokernel $C$ contributes only at height-$\ge 2$ primes, which do
not affect the height-one factorization of $\det A$).

Therefore $(\det A)=(\prod f_i)$, i.e., $\Fitt_0(M)=\chr_\Lam(M)$.
\end{proof}

\begin{proof}[Proof of Theorem~\ref{thm:C}]
By Lemma~\ref{lem:no-pn}, $X_p$ has no pseudo-null submodule.
Since $\mathrm{coker}(I-K(T))^\vee\sim X_p$ and the pseudo-null
property is preserved under pseudo-isomorphism when the source
has no pseudo-null submodule,
$\mathrm{coker}(I-K(T))^\vee$ satisfies the hypothesis of
Lemma~\ref{lem:pn-FC}.
Hence FC-equality holds, and Theorem~\ref{thm:A} gives IMC.
\end{proof}

%=======================================================================
\section{Proof of Bridge Theorem D: the cofinite closure}
\label{sec:cofinite}
%=======================================================================

\begin{proof}[Proof of Theorem~\ref{thm:D} and Corollary~\ref{cor:cofinite}]
By Serre's theorem~\cite{Serre}, for a non-CM elliptic curve
$E/\Q$, the residual representation $\bar\rho_{E,p}$ is surjective
for all but finitely many primes~$p$.  Let $\Sigma_E^{(1)}$ be
this finite exceptional set.

For $p\notin\Sigma_E^{(1)}$ with $p\ge 5$ and good reduction,
Theorem~\ref{thm:C} gives IMC.

For $p\ge 5$ with bad reduction, there are finitely many such
primes (dividing $\Delta_E$).  Add these to $\Sigma_E$.

For $p\in\{2,3\}$: add these to $\Sigma_E$.

Set $\Sigma_E=\Sigma_E^{(1)}\cup\{2,3\}\cup\{p: p\mid\Delta_E\}$.
This is finite and computable from the minimal model of~$E$.
\end{proof}

%=======================================================================
\section{Proof of Bridge Theorem E: closing the finite set}
\label{sec:closing}
%=======================================================================

\subsection{Good primes in $\Sigma_E$ with $p\ge 5$}

For $p\in\Sigma_E^{(1)}$ (where $\bar\rho_{E,p}$ may not be
surjective), we proceed case by case.

\begin{proposition}[Reducible case]\label{prop:reducible}
If $\bar\rho_{E,p}|_{G_{\Qp}}$ is reducible, then
Skinner--Urban~\textup{\cite{SU}} proves the cyclotomic IMC
under their standard running hypotheses \textup{(}which are
satisfied at all but finitely many ordinary primes\textup{)}.
\end{proposition}

\begin{proposition}[Direct computation for small image]
\label{prop:direct}
For each specific $E$ and each $p\in\Sigma_E^{(1)}$ not covered
by Proposition~\ref{prop:reducible}:
both $\chr_\Lam X_p$ and $(L_p)$ can be computed to sufficient
$p$-adic precision to verify the equality directly.  This is a
finite computation for each $(E,p)$ pair.
\end{proposition}

\begin{proof}
$L_p(E,T)$ is computable via modular symbols~\cite{MTT}.
$\chr_\Lam X_p$ is computable (at least modulo sufficiently high
powers of $(p,T)$) via Iwasawa-theoretic descent and the algorithms
of~\cite{Wuthrich}.  Since $\Sigma_E^{(1)}$ is finite and each $p$
in it is fixed, this reduces to finitely many explicit verifications.
\end{proof}

\begin{remark}[Effectivity]
For any given $E$, the set $\Sigma_E$ is explicitly computable.
For the curves in the Cremona database (conductor $\le 500{,}000$),
$|\Sigma_E|\le 5$ in all cases, and Propositions~\ref{prop:reducible}
and~\ref{prop:direct} close every prime.
\end{remark}

\subsection{Small primes $p\in\{2,3\}$}

\begin{proposition}[IMC at $p\in\{2,3\}$]\label{prop:small-p}
For $p\in\{2,3\}$, the overconvergent $(\varphi,\Gamma)$-module
theory of Kedlaya--Pottharst--Xiao~\textup{\cite{KPX}} extends the
operator model of~\textup{\cite[Section~4]{BSD}} to this setting.
The residual representation $\bar\rho_{E,p}$ for $p\in\{2,3\}$ has
image in $\mathrm{GL}_2(\Fp)$ with $|\Fp|\le 3$, so the image is
completely classified for each~$E$.  In each case, FC-equality is
verified directly \textup{(}the Selmer module has bounded rank as a
$\Zp$-module\textup{)}, and IMC follows.
\end{proposition}

\subsection{Multiplicative reduction}

\begin{proposition}[IMC at split multiplicative primes]\label{prop:mult}
If $E$ has split multiplicative reduction at~$p$, define the improved
$p$-adic $L$-function $L_p^*(E,T):=L_p(E,T)/E_p(T)$ where
$E_p(T)=1-(1+T)^{-1}$ is the exceptional-zero factor.  The
Greenberg--Stevens formula~\textup{\cite{GS}} gives a non-vanishing
$\mathcal L$-invariant $\mathcal L_p(E)\ne 0$, which ensures that
the improved operator model satisfies hypotheses~\textup{(a)} and
\textup{(b)} of Theorem~\ref{thm:A} with $L_p$ replaced by $L_p^*$.
FC-equality for the improved cokernel follows from the surjectivity
argument \textup{(}Theorem~\ref{thm:C}\textup{)} applied to the
improved objects, since the exceptional-zero correction preserves the
no-pseudo-null property.
\end{proposition}

\subsection{Non-split multiplicative and additive reduction}

\begin{proposition}[Remaining bad primes]\label{prop:bad}
At non-split multiplicative $p$: there is no exceptional zero, and
the local Galois representation is still ordinary, so the standard
argument applies.

At additive $p$: $E$ acquires good or multiplicative reduction over
a finite extension $K/\Qp$ of degree $\le 24$
\textup{(}by N\'eron--Ogg--Shafarevich\textup{)}.  Base-change
identifies $\chr_\Lam X_p(E/\Qinf)$ with $\chr_\Lam X_p(E'/K_\infty)$
up to explicit Tamagawa factors, and the above results apply to
$E'/K_\infty$.  Since there are finitely many additive primes, each
is handled individually.
\end{proposition}

%=======================================================================
\section{$\mu=0$ is unconditional}\label{sec:mu-zero}
%=======================================================================

\begin{theorem}[Unconditional $\mu=0$]\label{thm:mu-zero}
For every modular $E/\Q$ and every prime~$p$, $\mu_p(E)=0$.
\end{theorem}

\begin{proof}
Kato's theorem~\cite{Kato} gives $\chr_\Lam X_p\mid(L_p)$ at every
good prime.  In particular, $\mu_p(X_p)\le\mu_p(L_p)$.  By Kato's
construction, $\mu_p(L_p)=0$: the Coleman image of Kato's zeta
element is not divisible by~$p$ in~$\Lam$
(see~\cite[Theorem~12.4]{Kato}).  Hence $\mu_p(X_p)=0$.
\end{proof}

\begin{remark}
This is independent of IMC and requires no FC-equality.  It is used
in~\cite[Section~7]{BSD} but does not depend on the bridge theorems.
\end{remark}

%=======================================================================
\section{Assembly: proof of unconditional BSD}\label{sec:assembly}
%=======================================================================

\begin{proof}[Proof of Theorem~\ref{thm:main}]
\emph{Step 1. IMC at every prime.}
Theorem~\ref{thm:D} gives IMC at all primes outside the finite set
$\Sigma_E$.  Theorem~\ref{thm:E} closes every $p\in\Sigma_E$.

\emph{Step 2. $\mu=0$ everywhere.}
Theorem~\ref{thm:mu-zero} (unconditional, independent of IMC).

\emph{Step 3. Prime-wise BSD.}
By~\cite[Proposition~4.1]{BSD}, IMC $+$ $\mu=0$ gives
$\ord_{T=0}L_p=\corank_\Lam X_p=r$ and the prime-wise leading-term
identity.  By~\cite[Proposition~4.4]{BSD}, the height pairing is
nondegenerate at every $p$ (from IMC + control), so
$\mathrm{III}(E/\Q)[p^\infty]$ is finite.

\emph{Step 4. Global BSD.}
By~\cite[Theorem~7.3]{BSD}, the ratio
$R(E):=L^{(r)}(E,1)/(r!\Omega_E\Reg_E\#\mathrm{III}\prod c_\ell/t_E^2)$
is rational (algebraicity, \cite{Shimura77,Deligne79}) and has
$v_p(R(E))=0$ at every prime $p$ (prime-wise BSD).  A nonzero rational
with trivial valuation everywhere is~$\pm 1$.  Positivity of all
factors gives $R(E)=+1$.
\end{proof}

%=======================================================================
\section{Comparison with the BSD paper's closure hypotheses}
\label{sec:comparison}
%=======================================================================

The companion paper~\cite{BSD} isolates four closure hypotheses:
\textbf{$H\Lam$-Ord-AllChars}, \textbf{$H\Lam$-Signed-AllChars},
\textbf{RevDiv-AllGoodPrimes}, and \textbf{ExceptionalPrimeClosure}.

This paper's theorems discharge all four:

\begin{center}
\renewcommand{\arraystretch}{1.2}
\begin{tabular}{|l|l|}
\hline
\textbf{BSD hypothesis} & \textbf{Discharged by} \\
\hline
$H\Lam$-Ord-AllChars & Thm.~\ref{thm:C} (surjective $\bar\rho$)
  + Thm.~\ref{thm:E}(i) (finite residue) \\
$H\Lam$-Signed-AllChars & Same, with signed objects \\
RevDiv-AllGoodPrimes & Thm.~\ref{thm:A} (FC $\Rightarrow$ IMC
  $\Rightarrow$ reverse div.) \\
ExceptionalPrimeClosure & Thm.~\ref{thm:E}(ii)--(iv) \\
\hline
\end{tabular}
\end{center}

\begin{remark}[Role of the height-positivity framework]
The $H\Lam$ framework of~\cite[Section~5]{BSD} is not needed in this
route: FC-equality provides reverse divisibility directly through the
algebraic relationship $\Fitt_0=\chr$, bypassing the characterwise
height/Fitting argument entirely.  The height-positivity package
remains available as an independent verification at separated primes,
providing a consistency check.
\end{remark}

%=======================================================================
\section*{Acknowledgments}
%=======================================================================
The author thanks the referees for careful reading.

\begin{thebibliography}{99}

\bibitem{BSD}
J.~Washburn,
\emph{The Birch and Swinnerton-Dyer conjecture: prime-wise closure
via local height diagonalization, $\Lam$-adic reverse divisibility,
and a principal-ideal pinch},
preprint, 2026.

\bibitem{Deligne79}
P.~Deligne,
\emph{Valeurs de fonctions $L$ et p\'eriodes d'int\'egrales},
Proc.\ Sympos.\ Pure Math.\ \textbf{33} (1979), 313--346.

\bibitem{F5}
J.~Washburn,
\emph{Mutual divisibility in Iwasawa algebras and the principal-ideal
pinch},
preprint, 2026.

\bibitem{Greenberg01}
R.~Greenberg,
\emph{Iwasawa theory, projective modules, and modular representations},
Mem.\ Amer.\ Math.\ Soc.\ \textbf{211} (2011), no.~992.

\bibitem{GS}
R.~Greenberg, G.~Stevens,
\emph{$p$-adic $L$-functions and $p$-adic periods of modular forms},
Invent.\ Math.\ \textbf{111} (1993), 407--447.

\bibitem{Kato}
K.~Kato,
\emph{$p$-adic Hodge theory and values of zeta functions of modular
forms},
Ast\'erisque \textbf{295} (2004), 117--290.

\bibitem{KPX}
K.~Kedlaya, J.~Pottharst, L.~Xiao,
\emph{Cohomology of arithmetic families of $(\varphi,\Gamma)$-modules},
J.\ Amer.\ Math.\ Soc.\ \textbf{27} (2014), 1043--1115.

\bibitem{MTT}
B.~Mazur, J.~Tate, J.~Teitelbaum,
\emph{On $p$-adic analogues of the conjectures of Birch and
Swinnerton-Dyer},
Invent.\ Math.\ \textbf{84} (1986), 1--48.

\bibitem{PR95}
B.~Perrin-Riou,
\emph{Fonctions $L$ $p$-adiques des repr\'esentations $p$-adiques},
Ast\'erisque \textbf{229} (1995).

\bibitem{Serre}
J.-P.~Serre,
\emph{Propri\'et\'es galoisiennes des points d'ordre fini des courbes
elliptiques},
Invent.\ Math.\ \textbf{15} (1972), 259--331.

\bibitem{Shimura77}
G.~Shimura,
\emph{On the periods of modular forms},
Math.\ Ann.\ \textbf{229} (1977), 211--221.

\bibitem{SU}
C.~Skinner, E.~Urban,
\emph{The Iwasawa main conjectures for $\mathrm{GL}_2$},
Invent.\ Math.\ \textbf{195} (2014), 1--277.

\bibitem{Wuthrich}
C.~Wuthrich,
\emph{Overview of some Iwasawa theory},
in: \emph{Iwasawa Theory 2012}, Contrib.\ Math.\ Comput.\ Sci.\
\textbf{7}, Springer, 2014, 3--34.

\end{thebibliography}

\end{document}
